% Part: normal-modal-logic
% Chapter: tableaux
% Section: countermodels

\documentclass[../../../include/open-logic-section]{subfiles}

\begin{document}

\olfileid{nml}{tab}{cou}

\olsection{\usetoken{P}{tableau} توں ردّی نمونے}

\begin{explain}
  تمامیت قضیے دا ثبوت وکھاؤندا اے کہ جے $\Entails !A$، تاں
  $\Proves !A$۔ ایہ~$!A$ لئی ردّی نمونہ بناؤن دا طریقہ وی دیندا اے،
  جے $\Entails/ A$۔~$\Log{K}$ دی صورت وچ ایہ طریقہ \emph{فیصلہ کرن
    دا طریقہ} اے۔ منّ لو $\Entails/ !A$۔ تاں
  \olref[cpl]{prop:complete-tableau} دے ثبوت توں مکمل !!{tableau}
  بناؤن دا اک طریقہ ملدا اے۔ ایہ طریقہ ہمیشہ مک جاندا اے۔
  ~$\Log{K}$ دے قضیاتی قاعدے صرف گھٹ پیچیدگی والے اگیتے والے
  !!{formula}s جوڑدے نیں، یعنی ہر نشان والے فارمولے $\sFmla{S}{!A}[\sigma]$
  لئی اک شاخ اُتے ہر قضیاتی قاعدہ اک واری لانا کافی اے۔ نویں اگیتے صرف
  \iftag{prvBox}{$\TRule{\False}{\Box}$}{}\iftag{notprvBox,notprvDiamond}{}{
    تے }\iftag{prvDiamond}{$\TRule{\True}{\Diamond}$}{}
  \iftag{notprvBox,notprvDiamond}{قاعدے}{قاعدیاں} توں ملدے نیں۔
  ایہ وی اک واری لانے کافی نیں (تے اکّو نواں اگیتا دیندے نیں)۔
  \iftag{prvBox}{$\TRule{\True}{\Box}$}{}\iftag{notprvBox,notprvDiamond}{}{
    تے }\iftag{prvDiamond}{$\TRule{\False}{\Diamond}$}{}
  \iftag{notprvBox,notprvDiamond}{قاعدہ}{قاعدے} اک توں ودھ وار لانے پے
  سکدے نیں، پر ہر اگیتے لئی اکّو واری۔ نویں اگیتے محدود گنتی وچ ای
  بندے نیں۔ سو محدود پڑاواں پچھوں ایس بناوٹ توں یا بند شاخ ملدی اے
  یا مکمل شاخ۔

جدوں کھلی مکمل شاخ والا تابلو بن جائے، تاں
  \olref[cpl]{thm:tableau-completeness} دا ثبوت سانوں اوہ عین نمونہ
  دیندا اے جیہڑا اگیتے والے !!{formula}s دے اصل سیٹ نوں پورا کردا اے۔
  سو جے $\Gamma \Entails !A$، تاں بند !!{tableau} موجود اے تے
  $\Gamma \Proves !A$۔ تے جے بند !!{tableau} نوں ٹھیک طریقے نال لبھدیاں
  اک «مکمل» !!{tableau} ملے، تاں سانوں $\Gamma \Entails/ !A$ دا پتا
  لگے گا تے اسیں ردّی نمونہ وی بنا سکاں گے۔
\end{explain}

\iftag{prvBox}{
\begin{ex}
  سانوں پتا اے کہ $\Proves/ \Box(p \lor q) \lif (\Box p \lor \Box
  q)$۔ تابلو بناؤن دا مڈھ ایویں اے:
  \begin{oltableau}
    [\pFmla{\False}{\Box(p \lor q) \lif (\Box p \lor \Box q)}{1},
      just = \TAss, checked
      [\pFmla{\True}{\Box(p \lor q)}{1},
        just = {\TRule{\False}{\lif}[1]}, 
        [\pFmla{\False}{\Box p \lor \Box q}{1},
          just = {\TRule{\False}{\lif}[1]}, checked
          [\pFmla{\False}{\Box p}{1},
            just = {\TRule{\False}{\lor}[3]}, checked
            [\pFmla{\False}{\Box q}{1},
              just = {\TRule{\False}{\lor}[3]}, checked
              [\pFmla{\False}{p}{1.1}, 
                just = {\TRule{\False}{\Box}[4]}, checked
                [\pFmla{\False}{q}{1.2}, 
                  just = {\TRule{\False}{\Box}[5]}, checked
                ]
              ]
            ]
          ]
        ]
      ]
    ]
  \end{oltableau}
  !!{tableau} ہجے پورا نہیں ہویا۔ اگلے پڑا وچ اسیں اوہ اکّو سطر ویکھدے
  آں جیہدے اُتے صحیح دا نشان نہیں، یعنی اگیتے والا !!{formula}
  $\sFmla{\True}{\Box(p \lor q)}[1]$ سطر~$2$ اُتے۔ بند تابلو بناؤن لئی
  $\TRule{\True}{\Box}$ دا قاعدہ شاخ اُتے ورتے ہوئے ہر اگیتے لئی لانا
  اے، یعنی $1.1$ تے $1.2$ دونہاں لئی:
  \begin{oltableau}
    [\pFmla{\False}{\Box(p \lor q) \lif (\Box p \lor \Box q)}{1},
      just = \TAss, checked
      [\pFmla{\True}{\Box(p \lor q)}{1},
        just = {\TRule{\False}{\lif}[1]}, 
        [\pFmla{\False}{\Box p \lor \Box q}{1},
          just = {\TRule{\False}{\lif}[1]}, checked
          [\pFmla{\False}{\Box p}{1},
            just = {\TRule{\False}{\lor}[3]}, checked
            [\pFmla{\False}{\Box q}{1},
              just = {\TRule{\False}{\lor}[3]}, checked
              [\pFmla{\False}{p}{1.1}, 
                just = {\TRule{\False}{\Box}[4]}, checked
                [\pFmla{\False}{q}{1.2}, 
                  just = {\TRule{\False}{\Box}[5]}, checked
                  [\pFmla{\True}{p \lor q}{1.1}, 
                    just = {\TRule{\True}{\Box}[2]}
                    [\pFmla{\True}{p \lor q}{1.2}, 
                      just = {\TRule{\True}{\Box}[2]}
                    ]
                  ]
                ]
              ]
            ]
          ]
        ]
      ]
    ]
  \end{oltableau}
  ہُن سطر 2، 8 تے 9 اُتے صحیح دے نشان نہیں۔ پر کوئی نواں اگیتا نہیں
  جڑیا، ایس لئی اسیں $\TRule{\True}{\lor}$ سطر 8 تے 9 اُتے، ملن والیاں
  ساریاں شاخاں وچ (جنا چر اوہ بند نہیں ہوندیاں)، لاؤندے آں:
  \begin{oltableau}
    [\pFmla{\False}{\Box(p \lor q) \lif (\Box p \lor \Box q)}{1},
      just = \TAss, checked
      [\pFmla{\True}{\Box(p \lor q)}{1},
        just = {\TRule{\False}{\lif}[1]}, checked
        [\pFmla{\False}{\Box p \lor \Box q}{1},
          just = {\TRule{\False}{\lif}[1]}, checked
          [\pFmla{\False}{\Box p}{1},
            just = {\TRule{\False}{\lor}[3]}, checked
            [\pFmla{\False}{\Box q}{1},
              just = {\TRule{\False}{\lor}[3]}, checked
              [\pFmla{\False}{p}{1.1}, 
                just = {\TRule{\False}{\Box}[4]}, checked
                [\pFmla{\False}{q}{1.2}, 
                  just = {\TRule{\False}{\Box}[5]}, checked
                  [\pFmla{\True}{p \lor q}{1.1}, 
                    just = {\TRule{\True}{\Box}[2]}, checked
                    [\pFmla{\True}{p \lor q}{1.2}, 
                      just = {\TRule{\True}{\Box}[2]}, checked
                      [\pFmla{\True}{p}{1.1},
                        just = {\TRule{\True}{\lor}[8]}, checked, close
                      ]
                      [\pFmla{\True}{q}{1.1},
                        just = {\TRule{\True}{\lor}[8]}, checked
                        [\pFmla{\True}{p}{1.2},
                          just = {\TRule{\True}{\lor}[9]}, checked]
                        [\pFmla{\True}{q}{1.2},
                          just = {\TRule{\True}{\lor}[9]}, checked, close]
                      ]
                    ]
                  ]
                ]
              ]
            ]
          ]
        ]
      ]
    ]
  \end{oltableau}
  اک کھلی شاخ بچدی اے، تے اوہ مکمل اے۔ اوس توں اسیں نمونہ متعین کردے آں۔
  جہان $W = \{1, 1.1, 1.2\}$ نیں (کھلی شاخ اُتے آؤن والے ای اگیتے)،
  رسائی دا تعلق $R = \{\tuple{1, 1.1}, \tuple{1, 1.2}\}$ اے، تے قدر بندی
  $V(p) = \{1.2\}$ اے (کیوں جے سطر 11 اُتے $\sFmla{\True}{p}[1.2]$ اے)
  تے $V(q) = \{1.1\}$ اے (کیوں جے سطر 10 اُتے $\sFmla{\True}{q}[1.1]$ اے)۔
  نمونے دی تصویر \olref{fig:counter-Box} وچ اے، تے تسی پرکھ سکدے او کہ ایہ
  $\Box(p \lor q) \lif (\Box p \lor \Box q)$ دا ردّی نمونہ اے۔
  \begin{figure}
  \begin{center}
    \begin{tikzpicture}[modal]
      \node[world] (w1) [label={[align=right]right:\mFalse{p}\\ \mFalse{q}}]{$1$}; 
      \node[world] (w2) [label={[align=right]right:\mFalse{p}\\ \mTrue{q}},
        above left=of w1]{$1.1$}; 
      \node[world] (w3) [label={[align=right]right:\mTrue{p}\\ \mFalse{q}},
        above right=of w1] {$1.2$};
      \draw[->] (w1) to (w2);
      \draw[->] (w1) to (w3);
    \end{tikzpicture}
  \end{center}
  \caption{$\Box(p \lor q) \lif (\Box p \lor \Box
    q)$ دا اک ردّی نمونہ۔}
  \ollabel{fig:counter-Box}
\end{figure}
\end{ex}
}
{
\begin{ex}
  سانوں پتا اے کہ $\Proves/ (\Diamond p \land \Diamond q) \lif \Diamond(p
  \land q)$۔ تابلو بناؤن دا مڈھ ایویں اے:
  \begin{oltableau}
    [\pFmla{\False}{(\Diamond p \land \Diamond q) \lif \Diamond(p \land q)}{1},
      just = \TAss, checked
      [\pFmla{\True}{\Diamond p \land \Diamond q}{1},
        just = {\TRule{\False}{\lif}[1]}, checked
        [\pFmla{\False}{\Diamond(p \land q)}{1},
          just = {\TRule{\False}{\lif}[1]}
          [\pFmla{\True}{\Diamond p}{1},
            just = {\TRule{\True}{\land}[2]}, checked
            [\pFmla{\True}{\Diamond q}{1},
              just = {\TRule{\True}{\land}[2]}, checked
              [\pFmla{\True}{p}{1.1}, 
                just = {\TRule{\True}{\Diamond}[4]}, checked
                [\pFmla{\True}{q}{1.2}, 
                  just = {\TRule{\True}{\Diamond}[5]}, checked
                ]
              ]
            ]
          ]
        ]
      ]
    ]
  \end{oltableau}
  !!{tableau} ہجے پورا نہیں ہویا۔ اگلے پڑا وچ اسیں اوہ اکّو سطر ویکھدے
  آں جیہدے اُتے صحیح دا نشان نہیں، یعنی اگیتے والا !!{formula}
  $\sFmla{\True}{\Diamond(p \land q)}[1]$ سطر~$3$ اُتے۔ بند تابلو بناؤن لئی
  $\TRule{\True}{\Diamond}$ دا قاعدہ شاخ اُتے ورتے ہوئے ہر اگیتے لئی لانا
  اے، یعنی $1.1$ تے $1.2$ دونہاں لئی:
  \begin{oltableau}
    [\pFmla{\False}{\Diamond(p \land q) \lif (\Diamond p \land \Diamond q)}{1},
      just = \TAss, checked
      [\pFmla{\True}{\Diamond p \land \Diamond q}{1},
        just = {\TRule{\False}{\lif}[1]}, checked
        [\pFmla{\False}{\Diamond (p \land q)}{1},
          just = {\TRule{\False}{\lif}[1]}
          [\pFmla{\True}{\Diamond p}{1},
            just = {\TRule{\True}{\land}[2]}, checked
            [\pFmla{\True}{\Diamond q}{1},
              just = {\TRule{\True}{\land}[2]}, checked
              [\pFmla{\True}{p}{1.1}, 
                just = {\TRule{\True}{\Diamond}[4]}, checked
                [\pFmla{\True}{q}{1.2}, 
                  just = {\TRule{\True}{\Diamond}[5]}, checked
                  [\pFmla{\False}{p \land q}{1.1}, 
                    just = {\TRule{\False}{\Diamond}[3]}
                    [\pFmla{\False}{p \land q}{1.2}, 
                      just = {\TRule{\False}{\Diamond}[3]}
                    ]
                  ]
                ]
              ]
            ]
          ]
        ]
      ]
    ]
  \end{oltableau}
  ہُن سطر 3، 8 تے 9 اُتے صحیح دے نشان نہیں۔ پر کوئی نواں اگیتا نہیں
  جڑیا، ایس لئی اسیں $\TRule{\False}{\land}$ سطر 8 تے 9 اُتے، ملن والیاں
  ساریاں شاخاں وچ (جنا چر اوہ بند نہیں ہوندیاں)، لاؤندے آں:
  \begin{oltableau}
    [\pFmla{\False}{(\Diamond p \land \Diamond q) \lif \Diamond(p \land q)}{1},
      just = \TAss, checked
      [\pFmla{\True}{\Diamond p \land \Diamond q}{1},
        just = {\TRule{\False}{\lif}[1]}, checked
        [\pFmla{\False}{\Diamond(p \land q)}{1},
          just = {\TRule{\False}{\lif}[1]}
          [\pFmla{\True}{\Diamond p}{1},
            just = {\TRule{\True}{\land}[2]}, checked
            [\pFmla{\True}{\Diamond q}{1},
              just = {\TRule{\True}{\land}[2]}, checked
              [\pFmla{\True}{p}{1.1}, 
                just = {\TRule{\True}{\Diamond}[4]}, checked
                [\pFmla{\True}{q}{1.2}, 
                  just = {\TRule{\True}{\Diamond}[5]}, checked
                  [\pFmla{\False}{p \land q}{1.1}, 
                    just = {\TRule{\False}{\Diamond}[3]}, checked
                    [\pFmla{\False}{p \land q}{1.2}, 
                      just = {\TRule{\False}{\Diamond}[3]}, checked
                      [\pFmla{\False}{p}{1.1},
                        just = {\TRule{\False}{\land}[8]}, checked, close
                      ]
                      [\pFmla{\False}{q}{1.1},
                        just = {\TRule{\False}{\land}[8]}, checked
                        [\pFmla{\False}{p}{1.2},
                          just = {\TRule{\False}{\land}[9]}, checked]
                        [\pFmla{\False}{q}{1.2},
                          just = {\TRule{\False}{\land}[9]}, checked, close]
                      ]
                    ]
                  ]
                ]
              ]
            ]
          ]
        ]
      ]
    ]
  \end{oltableau}
  اک کھلی شاخ بچدی اے، تے اوہ مکمل اے۔ اوس توں اسیں نمونہ متعین کردے آں۔
  جہان $W = \{1, 1.1, 1.2\}$ نیں (کھلی شاخ اُتے آؤن والے ای اگیتے)،
  رسائی دا تعلق $R = \{\tuple{1, 1.1}, \tuple{1, 1.2}\}$ اے، تے قدر بندی
  $V(p) = \{1.1\}$ اے (کیوں جے سطر 6 اُتے $\sFmla{\True}{p}[1.1]$ اے)
  تے $V(q) = \{1.2\}$ اے (کیوں جے سطر 7 اُتے $\sFmla{\True}{q}[1.1]$ اے)۔
  نمونے دی تصویر \olref{fig:counter-Diamond} وچ اے، تے تسی پرکھ سکدے او کہ ایہ
  $(\Diamond p \land \Diamond q) \lif \Diamond (p \land q)$ دا ردّی نمونہ اے۔
  \begin{figure}
  \begin{center}
    \begin{tikzpicture}[modal]
      \node[world] (w1) [label={[align=right]right:\mFalse{p}\\ \mFalse{q}}]{$1$}; 
      \node[world] (w2) [label={[align=right]right:\mTrue{p}\\ \mFalse{q}},
        above left=of w1]{$1.1$}; 
      \node[world] (w3) [label={[align=right]right:\mFalse{p}\\ \mTrue{q}},
        above right=of w1] {$1.2$};
      \draw[->] (w1) to (w2);
      \draw[->] (w1) to (w3);
    \end{tikzpicture}
  \end{center}
  \caption{$(\Diamond p \land \Diamond q) \lif
    \Diamond (p \land q)$ دا اک ردّی نمونہ۔}
  \ollabel{fig:counter-Diamond}
\end{figure}
\end{ex}
}

\end{document}
